\uuid{3ywf}
\exo7id{4720}
\titre{exo7 4720}
\auteur{quercia}
\organisation{exo7}
\datecreate{2010-03-16}
\isIndication{false}
\isCorrection{true}
\chapitre{Topologie}
\sousChapitre{Topologie de la droite réelle}
\cpge{mpsi-mp2i.structures-algebriques.b, mpsi-mp2i.nombres-reels-suites.b, pcsi.nombres-reels-suites.a, ptsi.nombres-reels-suites.a, mp-mpi.topologie-espaces-normes.d, pc.espaces-vectoriels-normes.d, psi.espaces-vectoriels-normes.d, pt.fonctions-deux-trois-variables.a, mp-mpi.structures-algebriques.a}
\module{Analyse}
\niveau{L2}
\difficulte{}

\contenu{
\question{
Soit $A =\{ n+p\sqrt2 \text{ tq } n,p \in \N,\ n+p\sqrt2 > 0,\ n^2-2p^2 = 1\}$.
Montrer que $A$ est un sous-groupe discret de $\R^{+*}$.
}
\reponse{
$n+p\sqrt2 > 1  \Rightarrow  n>0,\ p>0$, donc $A\cap]1,+\infty[$
         admet un plus petit {\'e}l{\'e}ment : $3+2\sqrt2$.
}
}
